\paragrafos{Δενδροδιάγραμμα και πίνακας}{Γ.3.2, «Δενδροδιάγραμμα»}

\begin{anagke}
Στο κυλικείο του σχολείου η Φανή και ο Ηλίας θέλουν να δουν όλα τα δυνατά γεύματα. Υπάρχουν δύο είδη σάντουιτς, με τυρί ή με γαλοπούλα, και τρία ροφήματα: νερό, χυμό ή γάλα. Ο Ηλίας αρχίζει να τα γράφει στη σειρά: «τυρί με νερό, γαλοπούλα με γάλα, τυρί με χυμό…». Σε λίγο μπερδεύεται: το «τυρί με γάλα» το έγραψε ήδη ή όχι;

Η Φανή προτείνει: «Να γράφουμε πρώτα το σάντουιτς και από κάθε σάντουιτς να ξεκινά ένα κλαδί για κάθε ρόφημα». Με αυτόν τον τρόπο δεν ξεχνάμε κανένα συνδυασμό και δεν γράφουμε κανένα δύο φορές.

Σε αυτή την παράγραφο μαθαίνουμε να βρίσκουμε τον δειγματικό χώρο πειραμάτων με περισσότερα στάδια, με δενδροδιάγραμμα και με πίνακα διπλής εισόδου, και να μετράμε τα αποτελέσματα με πολλαπλασιασμό.
\end{anagke}

\begin{lexi}{δενδροδιάγραμμα, κλαδί, πίνακας διπλής εισόδου}{ονόματα και σημασίες}
Το \lexiword{δενδροδιάγραμμα} είναι ένα διάγραμμα που μοιάζει με δέντρο που το έχουμε ξαπλώσει: έχει \emph{ρίζα}, \emph{κλαδιά} και \emph{φύλλα}. Ο \lexiword{πίνακας διπλής εισόδου} έχει δύο «εισόδους»: το πρώτο αποτέλεσμα στις γραμμές και το δεύτερο στις στήλες.
\end{lexi}

\begin{metafrasi}
Για κάθε επιλογή ξεκινά ένα κλαδί. & δενδροδιάγραμμα\\
Ένα μονοπάτι από τη ρίζα ως το φύλλο. & ένα αποτέλεσμα του πειράματος\\
Ένα κελί του πίνακα. & ένα αποτέλεσμα (γραμμή, στήλη)\\
Το πρώτο στάδιο έχει $3$ επιλογές και το δεύτερο $2$. & $3\cdot2=6$ αποτελέσματα\\
\end{metafrasi}

\tmima{Η θεωρία}

\begin{kanonas}{πώς φτιάχνω δενδροδιάγραμμα}
1. Ξεκινώ από μια αφετηρία (τη \emph{ρίζα}) και για κάθε δυνατό αποτέλεσμα του \textbf{πρώτου σταδίου} σχεδιάζω ένα κλαδί.\\
2. Από το τέλος κάθε κλαδιού σχεδιάζω νέα κλαδιά, ένα για κάθε δυνατό αποτέλεσμα του \textbf{επόμενου σταδίου}. Συνεχίζω για όλα τα στάδια.\\
3. Κάθε μονοπάτι από τη ρίζα ως το τέλος αντιστοιχεί σε ένα αποτέλεσμα. Τα γράφω στην τελευταία στήλη: αυτά είναι τα στοιχεία του δειγματικού χώρου.
\end{kanonas}

Παράδειγμα: οι δύο τύποι σάντουιτς, τυρί (Τ) και γαλοπούλα (Γ), και τα τρία ροφήματα, νερό (Ν), χυμός (Χ) και γάλα (Λ).

\begin{center}
\begin{tikzpicture}[scale=0.9]
\draw (2.00,-0.62)--(4.00,0.00);
\draw (2.00,-0.62)--(4.00,-0.62);
\draw (2.00,-0.62)--(4.00,-1.24);
\draw (2.00,-2.48)--(4.00,-1.86);
\draw (2.00,-2.48)--(4.00,-2.48);
\draw (2.00,-2.48)--(4.00,-3.10);
\draw (0.00,-1.55)--(2.00,-0.62);
\draw (0.00,-1.55)--(2.00,-2.48);
\node[tn] at (4.00,0.00) {Ν};
\node[anchor=west,font=\small\bfseries] at (4.45,0.00) {ΤΝ};
\node[tn] at (4.00,-0.62) {Χ};
\node[anchor=west,font=\small\bfseries] at (4.45,-0.62) {ΤΧ};
\node[tn] at (4.00,-1.24) {Λ};
\node[anchor=west,font=\small\bfseries] at (4.45,-1.24) {ΤΛ};
\node[tn] at (2.00,-0.62) {Τ};
\node[tn] at (4.00,-1.86) {Ν};
\node[anchor=west,font=\small\bfseries] at (4.45,-1.86) {ΓΝ};
\node[tn] at (4.00,-2.48) {Χ};
\node[anchor=west,font=\small\bfseries] at (4.45,-2.48) {ΓΧ};
\node[tn] at (4.00,-3.10) {Λ};
\node[anchor=west,font=\small\bfseries] at (4.45,-3.10) {ΓΛ};
\node[tn] at (2.00,-2.48) {Γ};
\node[tn] at (0,-1.55) {};
\node[font=\scriptsize,gray] at (2,0.8) {σάντουιτς};
\node[font=\scriptsize,gray] at (4,0.8) {ρόφημα};
\node[font=\scriptsize,gray,anchor=west] at (4.45,0.8) {αποτέλεσμα};
\end{tikzpicture}
\end{center}

Ο δειγματικός χώρος είναι $\{\mathrm{ΤΝ},\ \mathrm{ΤΧ},\ \mathrm{ΤΛ},\ \mathrm{ΓΝ},\ \mathrm{ΓΧ},\ \mathrm{ΓΛ}\}$ και έχει $6$ στοιχεία.

\begin{kanonas}{πλήθος αποτελεσμάτων}
Αν ένα πείραμα έχει στάδια και στο πρώτο υπάρχουν $m$ δυνατά αποτελέσματα, και μετά από κάθε αποτέλεσμα του πρώτου στο δεύτερο υπάρχουν $n$ δυνατά αποτελέσματα, τότε ο δειγματικός χώρος έχει $m\cdot n$ στοιχεία. Το ίδιο συνεχίζεται για περισσότερα στάδια: πολλαπλασιάζουμε το πλήθος των επιλογών κάθε σταδίου.
\end{kanonas}

Έτσι, αν ρίξουμε ένα νόμισμα $\nu$ φορές, τα αποτελέσματα είναι $2\cdot2\cdots2=2^\nu$. Για $3$ ρίψεις είναι $2^3=8$, για $4$ ρίψεις $2^4=16$.

\begin{kanonas}{πίνακας διπλής εισόδου}
Όταν το πείραμα έχει δύο στάδια, μπορούμε να βάλουμε τα αποτελέσματα του πρώτου σταδίου στις γραμμές και του δεύτερου στις στήλες. Κάθε κελί (γραμμή, στήλη) είναι ένα αποτέλεσμα του πειράματος. Έτσι, για δύο ρίψεις ενός ζαριού, ο πίνακας έχει $6\cdot6=36$ κελιά.
\end{kanonas}

\begin{prosoxi}
\begin{itemize}[leftmargin=1.2em,itemsep=0.3ex,topsep=2pt]
\item Από κάθε κλαδί του πρώτου σταδίου ξεκινούν κλαδιά για \emph{όλα} τα αποτελέσματα του δεύτερου σταδίου: μην ξεχάσεις κανένα.
\item Αν πάρω χωρίς επανατοποθέτηση, τα κλαδιά του δεύτερου σταδίου είναι λιγότερα, γιατί ό,τι βγήκε δεν μπορεί να ξαναβγεί.
\item Το πλήθος των αποτελεσμάτων προκύπτει με \emph{πολλαπλασιασμό}, όχι με πρόσθεση: για δύο ζάρια είναι $6\cdot6=36$ και όχι $6+6=12$.
\item Το $(2,5)$ και το $(5,2)$ είναι δύο διαφορετικά κελιά του πίνακα, δηλαδή διαφορετικά αποτελέσματα.
\end{itemize}
\end{prosoxi}

\tmima{Λυμένα παραδείγματα}

\paradeigma{Τρεις βολές}
\tixero Ο Νίκος κάνει τρεις βολές. Σε κάθε βολή ευστοχεί (Ε) ή αστοχεί (Α).\par
\tithelo Τον δειγματικό χώρο και πόσα αποτελέσματα έχουν ακριβώς δύο εύστοχες βολές.\par
\possindeo Φτιάχνω δενδροδιάγραμμα με τρία στάδια.\par

\begin{center}
\begin{tikzpicture}[scale=0.85]
\draw (3.20,-0.28)--(4.80,0.00);
\draw (3.20,-0.28)--(4.80,-0.55);
\draw (3.20,-1.38)--(4.80,-1.10);
\draw (3.20,-1.38)--(4.80,-1.65);
\draw (1.60,-0.82)--(3.20,-0.28);
\draw (1.60,-0.82)--(3.20,-1.38);
\draw (3.20,-2.48)--(4.80,-2.20);
\draw (3.20,-2.48)--(4.80,-2.75);
\draw (3.20,-3.58)--(4.80,-3.30);
\draw (3.20,-3.58)--(4.80,-3.85);
\draw (1.60,-3.03)--(3.20,-2.48);
\draw (1.60,-3.03)--(3.20,-3.58);
\draw (0.00,-1.93)--(1.60,-0.82);
\draw (0.00,-1.93)--(1.60,-3.03);
\node[tn] at (4.80,0.00) {Ε};
\node[anchor=west,font=\small\bfseries] at (5.25,0.00) {ΕΕΕ};
\node[tn] at (4.80,-0.55) {Α};
\node[anchor=west,font=\small\bfseries] at (5.25,-0.55) {ΕΕΑ};
\node[tn] at (3.20,-0.28) {Ε};
\node[tn] at (4.80,-1.10) {Ε};
\node[anchor=west,font=\small\bfseries] at (5.25,-1.10) {ΕΑΕ};
\node[tn] at (4.80,-1.65) {Α};
\node[anchor=west,font=\small\bfseries] at (5.25,-1.65) {ΕΑΑ};
\node[tn] at (3.20,-1.38) {Α};
\node[tn] at (1.60,-0.82) {Ε};
\node[tn] at (4.80,-2.20) {Ε};
\node[anchor=west,font=\small\bfseries] at (5.25,-2.20) {ΑΕΕ};
\node[tn] at (4.80,-2.75) {Α};
\node[anchor=west,font=\small\bfseries] at (5.25,-2.75) {ΑΕΑ};
\node[tn] at (3.20,-2.48) {Ε};
\node[tn] at (4.80,-3.30) {Ε};
\node[anchor=west,font=\small\bfseries] at (5.25,-3.30) {ΑΑΕ};
\node[tn] at (4.80,-3.85) {Α};
\node[anchor=west,font=\small\bfseries] at (5.25,-3.85) {ΑΑΑ};
\node[tn] at (3.20,-3.58) {Α};
\node[tn] at (1.60,-3.03) {Α};
\node[tn] at (0,-1.93) {};
\node[font=\scriptsize,gray] at (1.6,0.8) {1η βολή};
\node[font=\scriptsize,gray] at (3.2,0.8) {2η βολή};
\node[font=\scriptsize,gray] at (4.8,0.8) {3η βολή};
\end{tikzpicture}
\end{center}

\apantisi Ο δειγματικός χώρος έχει $2\cdot2\cdot2=8$ στοιχεία: $\mathrm{ΕΕΕ}$, $\mathrm{ΕΕΑ}$, $\mathrm{ΕΑΕ}$, $\mathrm{ΕΑΑ}$, $\mathrm{ΑΕΕ}$, $\mathrm{ΑΕΑ}$, $\mathrm{ΑΑΕ}$, $\mathrm{ΑΑΑ}$. Με ακριβώς δύο εύστοχες είναι τα $\mathrm{ΕΕΑ}$, $\mathrm{ΕΑΕ}$ και $\mathrm{ΑΕΕ}$, δηλαδή $3$ αποτελέσματα.\par
\elegxos Το πλήθος των φύλλων ($8$) συμφωνεί με το γινόμενο $2\cdot2\cdot2$.

\paradeigma{Ο κωδικός με τρία διαφορετικά ψηφία}
\tixero Η Πολυξένη φτιάχνει τριψήφιους κωδικούς με τα ψηφία $2$, $5$, $9$, χρησιμοποιώντας κάθε ψηφίο μόνο μία φορά.\par
\tithelo Όλους τους δυνατούς κωδικούς.\par
\possindeo Στο δενδροδιάγραμμα το πρώτο ψηφίο έχει $3$ επιλογές, το δεύτερο $2$ (ό,τι έμεινε) και το τρίτο $1$.\par

\begin{center}
\begin{tikzpicture}[scale=0.85]
\draw (3.20,0.00)--(4.80,0.00);
\draw (3.20,-0.55)--(4.80,-0.55);
\draw (1.60,-0.28)--(3.20,0.00);
\draw (1.60,-0.28)--(3.20,-0.55);
\draw (3.20,-1.10)--(4.80,-1.10);
\draw (3.20,-1.65)--(4.80,-1.65);
\draw (1.60,-1.38)--(3.20,-1.10);
\draw (1.60,-1.38)--(3.20,-1.65);
\draw (3.20,-2.20)--(4.80,-2.20);
\draw (3.20,-2.75)--(4.80,-2.75);
\draw (1.60,-2.48)--(3.20,-2.20);
\draw (1.60,-2.48)--(3.20,-2.75);
\draw (0.00,-1.38)--(1.60,-0.28);
\draw (0.00,-1.38)--(1.60,-1.38);
\draw (0.00,-1.38)--(1.60,-2.48);
\node[tn] at (4.80,0.00) {9};
\node[anchor=west,font=\small\bfseries] at (5.25,0.00) {259};
\node[tn] at (3.20,0.00) {5};
\node[tn] at (4.80,-0.55) {5};
\node[anchor=west,font=\small\bfseries] at (5.25,-0.55) {295};
\node[tn] at (3.20,-0.55) {9};
\node[tn] at (1.60,-0.28) {2};
\node[tn] at (4.80,-1.10) {9};
\node[anchor=west,font=\small\bfseries] at (5.25,-1.10) {529};
\node[tn] at (3.20,-1.10) {2};
\node[tn] at (4.80,-1.65) {2};
\node[anchor=west,font=\small\bfseries] at (5.25,-1.65) {592};
\node[tn] at (3.20,-1.65) {9};
\node[tn] at (1.60,-1.38) {5};
\node[tn] at (4.80,-2.20) {5};
\node[anchor=west,font=\small\bfseries] at (5.25,-2.20) {925};
\node[tn] at (3.20,-2.20) {2};
\node[tn] at (4.80,-2.75) {2};
\node[anchor=west,font=\small\bfseries] at (5.25,-2.75) {952};
\node[tn] at (3.20,-2.75) {5};
\node[tn] at (1.60,-2.48) {9};
\node[tn] at (0,-1.38) {};
\node[font=\scriptsize,gray] at (1.6,0.8) {1ο ψηφίο};
\node[font=\scriptsize,gray] at (3.2,0.8) {2ο ψηφίο};
\node[font=\scriptsize,gray] at (4.8,0.8) {3ο ψηφίο};
\end{tikzpicture}
\end{center}

\apantisi Οι κωδικοί είναι $259$, $295$, $529$, $592$, $925$, $952$. Συνολικά $6$, και $3\cdot2\cdot1=6$.\par
\elegxos Όλα τα ψηφία σε κάθε κωδικό είναι διαφορετικά και κανένας κωδικός δεν επαναλαμβάνεται.

\paradeigma{Δύο ζάρια, πίνακας διπλής εισόδου}
\tixero Ρίχνουμε δύο ζάρια, ένα κόκκινο και ένα μπλε, και προσθέτουμε τα δύο αποτελέσματα.\par
\tithelo Να δούμε τα αθροίσματα και πόσα ζευγάρια δίνουν άθροισμα $7$.\par
\possindeo Φτιάχνω πίνακα με τα αποτελέσματα του κόκκινου ζαριού στις γραμμές και του μπλε στις στήλες. Σε κάθε κελί γράφω το άθροισμα.\par

\begin{center}
\renewcommand{\arraystretch}{1.15}
\begin{tabular}{c|cccccc}
 & $1$ & $2$ & $3$ & $4$ & $5$ & $6$\\\hline
$1$ & $2$ & $3$ & $4$ & $5$ & $6$ & $7$\\
$2$ & $3$ & $4$ & $5$ & $6$ & $7$ & $8$\\
$3$ & $4$ & $5$ & $6$ & $7$ & $8$ & $9$\\
$4$ & $5$ & $6$ & $7$ & $8$ & $9$ & $10$\\
$5$ & $6$ & $7$ & $8$ & $9$ & $10$ & $11$\\
$6$ & $7$ & $8$ & $9$ & $10$ & $11$ & $12$\\
\end{tabular}
\end{center}

\apantisi Τα ζευγάρια είναι $6\cdot6=36$, ενώ τα διαφορετικά αθροίσματα είναι μόνο $11$ (από $2$ έως $12$). Άθροισμα $7$ δίνουν τα ζευγάρια $(1,6)$, $(2,5)$, $(3,4)$, $(4,3)$, $(5,2)$, $(6,1)$, δηλαδή $6$ κελιά.\par
\elegxos Τα $36$ κελιά είναι τα αποτελέσματα του πειράματος. Τα $11$ αθροίσματα είναι τιμές, που εμφανίζονται άλλες φορές λίγες και άλλες πολλές.

\paradeigma{Μέτρηση με πολλαπλασιασμό}
\tixero Η κλειδαριά ενός ποδηλάτου έχει τρεις τροχούς και σε καθέναν υπάρχουν τα ψηφία από $0$ έως $9$.\par
\tithelo Πόσους κωδικούς μπορεί να σχηματίσει;\par
\possindeo Κάθε τροχός έχει $10$ επιλογές, και οι τροχοί είναι τρία στάδια.\par
\apantisi $10\cdot10\cdot10=1000$ κωδικοί (από $000$ ως $999$).\par
\elegxos Οι αριθμοί από $000$ έως $999$ είναι $1000$.

\tmima{Ασκήσεις}

\zoni{Το κατάλαβα}{δέντρα και πίνακες}

\ex Ρίχνουμε ένα νόμισμα δύο φορές. Σχεδίασε δενδροδιάγραμμα και γράψε τον δειγματικό χώρο.

\ex Στον πίνακα φαίνεται το άθροισμα δύο ρίψεων ενός ζαριού. Συμπλήρωσέ τον. Μετά βρες πόσα ζευγάρια έχουν άθροισμα $8$ και πόσα διαφορετικά αθροίσματα υπάρχουν.\par
\begin{center}
\renewcommand{\arraystretch}{1.15}
\begin{tabular}{c|cccccc}
 & $1$ & $2$ & $3$ & $4$ & $5$ & $6$\\\hline
$1$ & $2$ & $3$ & $4$ & $5$ & $6$ & $7$\\
$2$ & $3$ & $\ldots$ & & & & \\
$3$ & $4$ & & & & & \\
$4$ & $5$ & & & & & \\
$5$ & $6$ & & & & & \\
$6$ & $7$ & & & & & \\
\end{tabular}
\end{center}

\ex Πόσα στοιχεία έχει ο δειγματικός χώρος όταν: α) ρίχνουμε ένα νόμισμα $3$ φορές, β) ρίχνουμε ένα ζάρι και ένα νόμισμα, γ) ρίχνουμε ένα ζάρι $2$ φορές, δ) ρίχνουμε ένα νόμισμα $4$ φορές;

\ex Μια πιτσαρία έχει δύο μεγέθη (μικρή, μεγάλη) και τρεις γεύσεις (μαργαρίτα, ζαμπόν, μανιτάρια). Σχεδίασε δενδροδιάγραμμα και γράψε όλες τις δυνατές επιλογές.

\ex Σωστό ή λάθος; α) Κάθε μονοπάτι ενός δενδροδιαγράμματος αντιστοιχεί σε ένα αποτέλεσμα. β) Σε τρεις ρίψεις νομίσματος υπάρχουν $6$ αποτελέσματα. γ) Σε έναν πίνακα διπλής εισόδου κάθε κελί είναι ένα αποτέλεσμα. δ) Δύο ρίψεις ενός ζαριού έχουν $12$ αποτελέσματα. ε) Ένα πείραμα με δύο στάδια, με $4$ επιλογές στο πρώτο και $5$ στο δεύτερο, έχει $20$ αποτελέσματα.

\ex Η Αγγελική έχει $4$ μπλούζες και $3$ φούστες. Πόσους διαφορετικούς συνδυασμούς μπορεί να φορέσει;

\zoni{Το χρησιμοποιώ}{γεύματα, κλειδαριές, αγώνες}

\ex \textbf{Οι καραμέλες.} Μια σακούλα έχει τρεις καραμέλες: μέντα (Μ), κανέλα (Κ) και μέλι (Λ). Βγάζουμε δύο, τη μία μετά την άλλη, χωρίς να ξαναβάλουμε την πρώτη. Σχεδίασε δενδροδιάγραμμα και γράψε τον δειγματικό χώρο.

\ex \textbf{Το πρωινό στο χωριό.} Η θεία Μαρία σερβίρει πρωινό: ψωμί (λευκό ή σκούρο), γέμιση (τυρί, αυγό ή ζαμπόν) και ρόφημα (γάλα ή χυμός). Σχεδίασε δενδροδιάγραμμα. Πόσα διαφορετικά πρωινά υπάρχουν;

\ex \textbf{Η κλειδαριά.} Η κλειδαριά ενός ντουλαπιού έχει τέσσερις τροχούς με ψηφία $0$ έως $9$. α) Πόσοι είναι οι δυνατοί κωδικοί; β) Ένας κωδικός με τρία ψηφία από το $1$ ως το $5$ πρέπει να ξεκινά από $1$ και να έχει όλα τα ψηφία διαφορετικά. Πόσους τέτοιους κωδικούς μπορεί να φτιάξει κανείς;

\ex \textbf{Η διαφορά.} Ο Ηλίας και η Ρήνη ρίχνουν από ένα ζάρι ο καθένας και υπολογίζουν τη διαφορά, μεγαλύτερο μείον μικρότερο (αν είναι ίσοι, η διαφορά είναι $0$). Φτιάξε πίνακα διπλής εισόδου. Πόσα ζευγάρια δίνουν διαφορά $0$ και πόσα διαφορά $5$;

\ex \textbf{Οι βολές.} Ο Νίκος κάνει τρεις βολές. Με δενδροδιάγραμμα βρες πόσα αποτελέσματα έχουν ακριβώς μία εύστοχη βολή και πόσα τουλάχιστον δύο εύστοχες βολές.

\zoni{Το σκέφτομαι}{γιατί δουλεύει}

\ex Συμπλήρωσε τον πίνακα για το πλήθος των αποτελεσμάτων όταν ρίχνουμε ένα νόμισμα $\nu$ φορές, για $\nu=1,2,3,4,5$. Πώς γράφεται το πλήθος ως δύναμη; Πόσα αποτελέσματα έχουν $10$ ρίψεις;

\ex Ο Βασίλης λέει: «Στα δύο ζάρια το $3$ και $4$ είναι το ίδιο με το $4$ και $3$». Ο Ηλίας διαφωνεί. Αν το ένα ζάρι είναι κόκκινο και το άλλο μπλε, ποιος έχει δίκιο όταν μετράμε ζευγάρια; Και όταν μας ενδιαφέρει μόνο το άθροισμα;

\ex Με τα ψηφία $1,2,3,4$ φτιάχνουμε τριψήφιους αριθμούς με όλα τα ψηφία διαφορετικά. Πόσους; (Σκέψου πόσες επιλογές έχει το κάθε ψηφίο.)

\ex Ο Στέλιος λέει ότι για τις δύο ρίψεις νομίσματος ο δειγματικός χώρος είναι $\{\mathrm{ΚΚ},\ \mathrm{ΚΓ},\ \mathrm{ΓΓ}\}$. Τι λείπει; Πώς θα το είχε καταλάβει με δενδροδιάγραμμα;

\zoni{Για το σπίτι}{στο σπίτι και στο παιχνίδι}

\ex Φτιάξε ένα δικό σου μενού ή ντουλάπα ρούχων με τρία στάδια και σχεδίασε δενδροδιάγραμμα. Πόσες είναι οι δυνατές επιλογές;

\ex Στο παιχνίδι «πέτρα, ψαλίδι, χαρτί» η πέτρα νικά το ψαλίδι, το ψαλίδι νικά το χαρτί και το χαρτί νικά την πέτρα. Φτιάξε πίνακα διπλής εισόδου για δύο παίκτες. Πόσα αποτελέσματα έχει το παιχνίδι; Σε πόσα από αυτά υπάρχει ισοπαλία;
